% Part: sets-functions-relations
% Chapter: relations
% Section: reflections
% Hindi translation (hi-Deva-IN) of Open Logic commit 9620cc73f9c8e0ad003c514a5d3748f29611c4c0.
%
\documentclass[../../../include/open-logic-section]{subfiles}

\begin{document}

\olfileid[hi]{sfr}{rel}{ref}
\olsection{दार्शनिक विचार}

\olref[set]{sec} में हमने संबंधों को कुछ विशेष समुच्चयों के रूप में परिभाषित
किया था। यहाँ रुककर एक छोटा-सा दार्शनिक प्रश्न पूछना चाहिए: ऐसी परिभाषा क्या
\emph{काम} कर रही है? यह कहना अत्यंत संदेहास्पद होगा कि हमने तत्त्वमीमांसात्मक तादात्म्य
के कुछ तथ्यों की \emph{खोज} कर ली है; कि उदाहरण के लिए, $\Nat$ पर क्रम-संबंध
\emph{वही निकला} समुच्चय $R=  \Setabs{\tuple{n,m}}{n, m \in \Nat\text{ और }
n < m}$ जिसे हमने \olref[set]{sec} में परिभाषित किया था। ऐसा सोचने के विरुद्ध
तीन कारण हैं।

पहला: \olref[set][pai]{wienerkuratowski} में हमने $\tuple{a, b} =
\{\{a\}, \{a, b\}\}$ परिभाषित किया था। इसके बजाय परिभाषा $\lVert a,
b\rVert = \{\{b\}, \{a, b\}\} = \tuple{b,a}$ पर विचार कीजिए। जब $a \neq b$,
तब $\tuple{a, b} \neq \lVert a,b\rVert$। किंतु हम समान रूप से
$\lVert a,b\rVert$ को क्रमित युग्म की अपनी परिभाषा मान सकते थे, न कि
$\tuple{a,b}$ को। दोनों
परिभाषाएँ समान रूप से अच्छी तरह काम करतीं। अतः अब प्राकृत संख्याओं पर क्रम-संबंध
``होने'' के दो समान रूप से अच्छे प्रत्याशी हैं, अर्थात्:
\begin{align*}
		R &= \Setabs{\tuple{n,m}}{n, m \in \Nat \text{ और }n < m}\\
		S &= \Setabs{\lVert n,m\rVert}{n, m \in \Nat \text{ और }n < m}.
\end{align*}
चूँकि विस्तारिता के अनुसार $R \neq S$, इसलिए स्पष्ट है कि दोनों एक साथ
\emph{उसी} क्रम-संबंध के तद्रूप नहीं हो सकते जो~$\Nat$ पर है। लेकिन यह दावा
करना कि तथ्यतः $R$, न कि $S$ (या इसका उलटा), \emph{ही} क्रम-संबंध है, बिल्कुल
मनमाना और इसलिए कुछ लज्जाजनक होगा। (यह समुच्चय-सैद्धांतिक अपचयनवाद के विरुद्ध
उस तर्क का अत्यंत सरल उदाहरण है जिसे \citeyear{Benacerraf1965} में बेनासेराफ
ने प्रसिद्ध बनाया। हम इस तर्क पर कई बार लौटेंगे।)

दूसरा: यदि हम मानें कि \emph{प्रत्येक} संबंध को किसी समुच्चय से तद्रूप माना
जाना चाहिए, तो स्वयं समुच्चय-अवयवता का संबंध~$\in$ भी कोई विशेष समुच्चय होना
चाहिए। वास्तव में उसे $\Setabs{\tuple{x,y}}{x \in y}$ समुच्चय होना पड़ेगा।
लेकिन क्या यह समुच्चय अस्तित्व में है? रसेल के विरोधाभास को देखते हुए, इस
समुच्चय का अस्तित्व कोई मामूली दावा नहीं है। वास्तव में,
\oliflabeldef{cumul:::part}{\olref[cumul][][]{part} में हम समुच्चयों का जो
सिद्धांत विकसित करेंगे, वह इस समुच्चय के अस्तित्व को \emph{अस्वीकार}
करेगा।\footnote{आगे की बात पहले ही बता दें, तो कारण यह है। प्रति-विरोध के लिए
मान लें कि $I =
\Setabs{\tuple{x,y}}{x \in y}$ अस्तित्व में है। तब $\bigcup \bigcup I$
सार्वत्रिक समुच्चय है, जो
\olref[sfr][z][sep]{thm:NoUniversalSet} का खंडन करता है।}}{समुच्चय-सिद्धांत को
स्वयंसिद्ध सिद्धांत के रूप में कठोर ढंग से विकसित करना संभव है, और वह सिद्धांत
वास्तव में इस समुच्चय के अस्तित्व को अस्वीकार करेगा।}
इसलिए कुछ संबंधों को समुच्चयों के रूप में माना जा सके, तब भी समुच्चय-अवयवता
का संबंध एक विशेष मामला होगा।

तीसरा: संबंधों को समुच्चयों से ``तद्रूप'' मानते समय हमने कहा था कि हम
$Rxy$ को $\tuple{x,y} \in R$ के लिए लिख सकेंगे। यह तब तक ठीक है जब
अवयवता-संबंध ``$\in$'' को किसी \emph{विधेय} के रूप में माना जाए। किंतु यदि हम
मानें कि ``$\in$'' किसी खास तरह के समुच्चय का द्योतक है, तो अभिव्यक्ति
``$\tuple{x,y} \in R$'' में केवल समुच्चयों के द्योतक तीन एकवाची पद रह जाएँगे:
``$\tuple{x,y}$'', ``$\in$'', और ``$R$''। नामों की ऐसी सूची किसी प्रतिज्ञप्ति
को व्यक्त करने में इस निरर्थक शब्द-माला से अधिक समर्थ नहीं है: ``कप कलमदान
मेज़''। फिर वही बात है: कुछ संबंधों को समुच्चयों के रूप में माना जा सके, तब भी
समुच्चय-अवयवता का संबंध एक विशेष मामला होना चाहिए। (इसमें फ़्रेगे के
\emph{संकल्पना ‘घोड़ा’} विरोधाभास के एक सरल रूप और रसेल के विरुद्ध
विट्गेन्श्टाइन की कभी उठाई गई प्रसिद्ध आपत्ति को एक साथ रख दिया गया है।)

तो इससे हम कहाँ पहुँचते हैं? संख्याओं पर बने संबंधों को समुच्चय कहने में कुछ
भी \emph{गलत} नहीं है। बस हमें उस अभिप्राय को समझना होगा जिसमें यह बात कही गई
है। हम तत्त्वमीमांसात्मक तादात्म्य का कोई तथ्य नहीं कह रहे हैं। हम केवल यह
देख रहे हैं कि कुछ संदर्भों में हम (कुछ) संबंधों को कुछ समुच्चयों के रूप में
\emph{बरत} सकते हैं (और बरतेंगे)।

\end{document}
